\begin{afterword}
\summaryhead

The post separates two ideas. ``Falsifiability'' is a property of one hypothesis: it concentrates its probability, so that some outcomes would nearly refute it. ``Testability'' is a relation between two: some outcome is more probable under one than under the other. Decoherent quantum mechanics, the author's name for the many-worlds view, is falsifiable, since observations such as apples falling upward would refute it; so is quantum mechanics with a collapse postulate.

On testability, the author argues that Bayes's theorem is symmetric. If decoherence is untestable relative to collapse, collapse is untestable relative to decoherence, and the order in which the theories were invented should not matter. What counts against a theory is a useless complication, new or old, as a parable about angels behind Maxwell's equations shows. The many worlds are not such a complication: they follow from the tested wave equation applied everywhere, while the collapse postulate can be dropped. The post concludes that excluding decoherence as extra entities, untestable or lacking new predictions is an error of probability theory.

\responsehead

The post's distinction is sound and its mathematics is correct. A single hypothesis can be judged by how much it forbids; choosing between two needs an outcome they assign different probabilities; Bayes's theorem does not care which theory came first. Its account of Popper on psychoanalysis is accurate. Its final claim is modest: decoherence should not be excluded on these three grounds, and it stays ``in the game.''

The title's last word is the weak point. By the post's own definition, a test needs an outcome that decoherence and collapse predict differently. The post names none. It argues instead that untestability would count against both theories equally, and that the shared wave equation is testable. Such outcomes do exist in principle: a defender of many-worlds, Vaidman, writes that the claim that it cannot be told apart from an ideal collapse theory ``is not the case,'' and experiments have ruled out some physical collapse models. ``Collapse Postulates,'' placed two posts earlier in the book, makes this point.

Two steps are stated as settled that are in dispute: that to deny the worlds is ``necessarily'' to deny that the quantum laws hold everywhere, which Bohmian mechanics contests, and that decoherence is ``strictly simpler,'' although the post's own version adds back ``the Born probability rule.'' Both are assessed in the notes on ``Many Worlds, One Best Guess.''

The post treats the rival view on new predictions unfairly. Predictivism, the view that a successful prediction counts for more than a fitted result, is Popper's position. The post answers it as order-dependence in updating, but its Bayesian defenders, such as Maher, argue that a successful prediction is itself extra evidence about the theory, which the theorems showing that Bayesian updating does not depend on the order of evidence do not rule out.

A smaller point: the post suggests, with ``may be,'' that this is the first time scientists have said a successful theory contains a part that can be dropped. The ether, dropped as superfluous in 1905, is a close precedent, and one that supports the post's rule.

In short: a correct distinction and correct mathematics, which keep many-worlds in contention, under a title whose ``testable'' is argued by symmetry rather than shown, with two disputed steps, ``necessarily'' and ``strictly simpler,'' stated as settled.
\end{afterword}
